\section{AI智能体的不足}

尽管前文已经展示了 AI 智能体在资料归集、数据整理和规则化处理中的显著优势，但这并不意味着其能力可以无条件外推到所有业务环节。特别是在安全要求高、环境变化快、需要持续情境判断和即时协同处置的现场作业中，AI 智能体目前仍存在明显边界。

\subsection{现场巡检与生产作业并非纯信息处理任务}

笔者在 2025 年底于油库工作期间，曾听到关于“以机器人替代人工巡检”的讨论。这一设想在技术演示层面确实具有吸引力，但若据此推断当前 AI 或机器人系统已经能够直接替代油库巡检岗位，显然仍为时过早。以油库巡检为例，完整路线一次最快也需要约 \SI{30}{\minute}，行走距离超过\SI{2}{\kilo\meter}。在这一过程中，工作人员并不是简单地沿既定路线行走，而是需要持续结合气味、温度、声音、震动、仪表状态、地面积液、设备细微变化以及现场作业上下文进行综合判断，并在发现异常时立即与中控或相关岗位联动处置。

以宇树科技的相关产品资料为例，其官方说明已经展示出较强的移动与巡检潜力，例如可通过激光雷达建图、按自定义路网进行自主巡逻与导航，并实时显示状态、温度和报警等信息\cite{unitree_b2_app_2025}。这说明当前机器人已经能够承担一部分“按路线移动+采集部分状态信息”的工作。然而，这与真正意义上的现场安全判断仍存在明显差距。前者更接近预设路径下的移动执行与状态回传，后者则需要面向开放环境进行连续的多模态理解、异常识别、风险分级以及处置决策。

同样地，春晚或技术大会上展示的人形机器人、娱乐机器人虽然令人印象深刻，但其“惊艳”主要体现于动作控制、交互表现、导航能力或策略训练进展，而不能直接等同于已具备在高风险工业环境中独立分析现实情境并承担安全责任的能力。以 NVIDIA 在 GTC 2026 期间发布的机器人生态进展为例，其官方表述强调的是 physical AI、世界模型、仿真训练、导航复杂环境与机器人策略学习等能力建设。例如 Disney 的 Olaf 和 BDX Droids 被用于训练策略并实现复杂环境导航\cite{nvidia_physical_ai_2026}。这说明机器人正朝着更强的感知与行动能力演进，但从“可演示、可训练、可导航”走向“可在危化品场景中稳定替代人工判断与处置”，中间仍有很长距离。

我还想补充：至少在油库场景中，AI 智能体至少暂时无法替代卸油栈桥上依靠经验判断现场状态的卸油人员，也难以替代泵房内持续拿着对讲机、与中控实时沟通并即时控制阀门的老师傅们。这些岗位的核心能力并不是单一动作执行，而是把经验、责任、风险意识和现场变化整合为即时判断。换言之，越是依赖复杂情境感知和责任闭环的现场岗位，越不能简单用当前的智能体能力去类比替代。

\subsection{岗位替代并非线性发生，现实影响仍在缓慢显现}

围绕 AI 是否会迅速替代大量岗位，近年来外部讨论往往较为激进。但从现有研究看，现实中的劳动替代效应并未以一种简单、同步、立刻显现的方式发生。Anthropic 于 2026 年发布的研究指出，AI 对劳动市场的影响需要区分“理论上能够完成哪些任务”和“现实中实际覆盖了多少任务”。该研究发现，AI 的实际覆盖范围仍明显低于其理论能力边界；同时，自 2022 年末以来，并未观察到高暴露岗位出现系统性的失业率上升，只是对年轻劳动者进入高暴露岗位的招聘放缓存在一定迹象\cite{anthropic_labor_market_2026}。

\begin{figure}[htbp]
    \centering
    \includegraphics[width=0.82\linewidth]{anthropic_labor-market-impacts.jpg}
    \caption{Anthropic 关于 AI 劳动市场影响研究中的相关图示\cite{anthropic_labor_market_2026}}
    \label{fig:anthropic-labor-market}
\end{figure}

这类研究结论对基层岗位判断具有重要启发。一方面，AI 的确正在改变部分任务的完成方式，尤其是文书、分析、客服、数据录入和编码等高信息密度场景；但另一方面，真正的岗位替代并不只是“模型能不能做”，还取决于制度责任、现场风险、组织流程、复核要求和软硬件协同条件。也就是说，AI 更可能先改变岗位内部的任务结构，而不是立刻整体消灭某个岗位。

\subsection{AI Agent 自身仍处于不够成熟的阶段}

除了物理现场作业难以完全替代之外，AI agent 自身的发展阶段也决定了它尚不适合被无条件神化。以软件工程领域的严格评测为例，SWE-Bench Pro 专门针对更接近真实企业环境、长时程、多文件修改的软件工程任务进行测试。结果表明，即便是当前领先模型，在该基准上的通过率仍低于 \SI{25}{\percent}，说明现阶段的 agent 在长链条任务、复杂上下文管理和企业级真实问题求解方面，距离稳定、鲁棒、可无人监督部署的状态还有较大差距\cite{deng2025swebenchpro}。

从产业观察角度看，Gartner 在 2025 年的预测也给出了类似提醒：到 2027 年底，超过 \SI{40}{\percent} 的 agentic AI 项目将被取消，其原因主要包括成本上升、业务价值不清晰以及风险控制不足；同时，Gartner 还指出，当前许多 agentic AI 项目仍停留在早期实验或概念验证阶段，且相当部分受到市场热度驱动\cite{gartner2025_agentic_cancelled}。这一判断说明，当前 agent 热潮中既包含真实进步，也夹杂着大量概念放大和场景误配。

因此，若回到基层经营管理实践，AI 智能体更适合被理解为“在特定规则和边界内显著提升效率的工具”，而不是能够脱离监管、脱离复核、脱离现场责任独立完成全部工作的替代者。对资料整理、表格分析、初稿生成、规则匹配这类任务，AI agent 已经具有较高价值；但对长时程、多环节、强责任、强协同、强现场感知的任务，其不足同样必须被正视。只有同时看到能力与边界，AI 智能体才能真正以务实方式服务于经营管理，而不至于陷入“技术神话”。
